% Abstract

Multi-scale network centrality is essential for urban analysis---from pedestrian accessibility at 500m to regional transport planning at 20km---but computationally prohibitive for large networks. Uniform sampling reduces cost but fails at short distances where few nodes are reachable, forcing a choice between accuracy and efficiency. We show that \emph{effective sample size}, the product of network reachability and sampling probability, modified by a power-law correction for sampling probability, predicts centrality estimation accuracy. This insight enables adaptive per-distance sampling: at each distance threshold, sampling probability is calibrated independently based on local reachability to maintain consistent accuracy. Empirically, harmonic closeness and betweenness centrality have distinct sample size requirements that vary by distance heuristic: for shortest paths, harmonic closeness requires approximately \ratioHBShortest$\times$ more samples than betweenness, while for angular paths betweenness requires approximately \ratioBHAngular$\times$ more. Validated on synthetic networks (grid, dendritic, linear topologies) and real-world street networks (London, Madrid, and Phoenix), adaptive sampling achieves speedups that increase with distance threshold while maintaining consistent accuracy. The approach is implemented in the open-source \texttt{cityseer} library, enabling rapid iteration for planning applications where computational cost currently limits exploratory analysis.
